\chapter*{Resumen}\label{Chap:Resumen}
\addcontentsline{toc}{chapter}{Resumen}

Este trabajo aborda el control de posición del sistema de levitación magnética ECP-730 en su configuración atractiva, inestable en lazo abierto, en presencia del efecto atascamiento-deslizamiento que produce el contacto del imán con su guía. El modelo no lineal emplea la fuerza electromagnética del fabricante y los parámetros identificados experimentalmente en trabajos previos; el atascamiento-deslizamiento se representa con la fricción estática de Karnopp y una zona muerta en el actuador, simulada como una retención de voltaje. Sobre el modelo linealizado se diseñó un controlador lineal con doble acción integral (PII), y su desempeño se evaluó por simulación en lazo abierto; sobre el modelo lineal con atascamiento-deslizamiento; sobre el modelo no lineal sin y con este efecto, en regulación y en seguimiento de referencias senoidal y trapezoidal; y en la configuración estable del mismo dispositivo. El PII estabiliza el sistema en todas las pruebas sin compensar explícitamente la fricción ni la zona muerta, con error medio prácticamente nulo y, en la configuración inestable, sin alcanzar el límite superior del actuador. Sin embargo, en las simulaciones el levitante no llega a detenerse sobre la referencia: oscila en un ciclo límite de 0,02 a 0,03~cm, causado por la interacción entre la acción integral y la fricción estática. La comparación con un controlador proporcional-integral (PI) y un barrido de dos familias de controladores PI y PII muestran que, en las familias y condiciones ensayadas, la segunda acción integral no reduce ni la amplitud del ciclo límite ni la duración de los atascamientos; la amplitud depende sobre todo de la ganancia proporcional. Se encontró además que, con un voltaje máximo de 3,5~V, el sistema solo puede cambiar de posición de forma regulada a partir de unos $-3{,}6$~cm, lo que restringe el intervalo de operación.
